\section{Stability Analysis}

In this section, the stability of the proposed adaptive neural controller is established using the Lyapunov direct method. For simplicity of the theoretical analysis, the nominal closed-loop system matrix $A$ defined in \eqref{eq:AB} is assumed to be Hurwitz. Consequently, for any given symmetric positive-definite matrix $Q\in\mathbb{R}^{2n\times2n}$, there exists a unique symmetric positive-definite matrix $P\in\mathbb{R}^{2n\times2n}$ satisfying the Lyapunov equation

\begin{equation}
	A^{T}P+PA=-Q,
	\label{eq:lyapunov}
\end{equation}

where $Q=Q^{T}>0$ and $P=P^{T}>0$.

%%%%%%%%%%%%%%%%%%%%%%%%%%%%%%%%%%%%%%%%%%%%%%%%%%%%%%%%%%%%%%

\subsection{Lyapunov Stability Analysis}

Consider the following composite Lyapunov candidate function

\begin{equation}
	V
	=
	\frac12
	X^{T}PX
	+
	\frac{1}{2\gamma}
	\operatorname{tr}
	\left(
	\tilde W^{T}\tilde W
	\right),
	\label{eq:V}
\end{equation}

where $\gamma>0$ denotes the adaptation gain.

Differentiating \eqref{eq:V} with respect to time yields

\begin{equation}
	\dot V
	=
	\frac12
	\dot X^{T}PX
	+
	\frac12
	X^{T}P\dot X
	+
	\frac1\gamma
	\operatorname{tr}
	\left(
	\dot{\tilde W}^{T}\tilde W
	\right).
	\label{eq:Vdot1}
\end{equation}

Since the ideal weight matrix $W^{*}$ is constant, one has

\begin{equation}
	\dot{\tilde W}
	=
	-
	\dot{\hat W}.
	\label{eq:weightdot}
\end{equation}

Substituting the closed-loop system \eqref{eq:ss}

\begin{equation}
	\dot X
	=
	AX
	+
	B
	\left(
	\tilde W^{T}h(x)
	+
	\varepsilon
	\right)
\end{equation}

into \eqref{eq:Vdot1} gives

\begin{align}
	\dot V
	=&
	\frac12
	\left[
	AX
	+
	B
	\left(
	\tilde W^{T}h
	+
	\varepsilon
	\right)
	\right]^{T}
	PX
	\nonumber\\
	&
	+
	\frac12
	X^{T}
	P
	\left[
	AX
	+
	B
	\left(
	\tilde W^{T}h
	+
	\varepsilon
	\right)
	\right]
	\nonumber\\
	&
	-
	\frac1\gamma
	\operatorname{tr}
	\left(
	\dot{\hat W}^{T}\tilde W
	\right).
	\label{eq:Vdot2}
\end{align}

Expanding \eqref{eq:Vdot2} yields

\begin{align}
	\dot V
	=&
	\frac12
	X^{T}
	(A^{T}P+PA)
	X
	\nonumber\\
	&
	+
	X^{T}PB
	\tilde W^{T}h
	+
	X^{T}PB\varepsilon
	\nonumber\\
	&
	-
	\frac1\gamma
	\operatorname{tr}
	\left(
	\dot{\hat W}^{T}\tilde W
	\right).
	\label{eq:Vdot3}
\end{align}

Using the Lyapunov equation \eqref{eq:lyapunov} $A^{T}P+PA=-Q$, one obtains

\begin{equation}
	\dot V
	=
	-
	\frac12
	X^{T}QX
	+
	X^{T}PB
	\tilde W^{T}h
	+
	X^{T}PB\varepsilon
	-
	\frac1\gamma
	\operatorname{tr}
	\left(
	\dot{\hat W}^{T}\tilde W
	\right).
	\label{eq:Vdot4}
\end{equation}

By exploiting the cyclic property of the trace operator,

\begin{equation}
	X^{T}PB\tilde W^{T}h
	=
	\operatorname{tr}
	\left(
	\tilde W^{T}
	h
	X^{T}
	PB
	\right),
\end{equation}

the derivative of the Lyapunov function can be rewritten as

\begin{align}
	\dot V
	=&
	-
	\frac12
	X^{T}QX
	+
	X^{T}PB\varepsilon
	\nonumber\\
	&
	+
	\frac1\gamma
	\operatorname{tr}
	\left[
	\tilde W^{T}
	\left(
	\gamma
	h
	X^{T}
	PB
	-
	\dot{\hat W}
	\right)
	\right].
	\label{eq:Vdot5}
\end{align}

\subsection{Adaptive Law with $\sigma$-Modification}

To eliminate the neural-network parameter estimation term in \eqref{eq:Vdot4}, the following trace identity is first employed:

\begin{equation}
	X^{T}PB\tilde{W}^{T}h(x)
	=
	\operatorname{tr}
	\left[
	\tilde{W}^{T}
	h(x)
	X^{T}PB
	\right].
	\label{eq:trace_identity}
\end{equation}

Substituting \eqref{eq:trace_identity} into \eqref{eq:Vdot4} yields

\begin{equation}
	\dot{V}
	=
	-\frac{1}{2}
	X^{T}QX
	+
	X^{T}PB\varepsilon
	+
	\frac{1}{\gamma}
	\operatorname{tr}
	\left[
	\tilde{W}^{T}
	\left(
	\gamma
	h(X)
	X^{T}PB
	-
	\dot{\hat{W}}
	\right)
	\right].
	\label{eq:adaptive1}
\end{equation}

Accordingly, the adaptive law is selected as

\begin{equation}
	\dot{\hat{W}}
	=
	\gamma
	h(X)
	X^{T}
	PB,
	\label{eq:adaptive2}
\end{equation}

which completely cancels the parameter estimation term in the Lyapunov derivative. Consequently,

\begin{equation}
	\dot{V}
	=
	-\frac{1}{2}
	X^{T}QX
	+
	X^{T}PB\varepsilon.
	\label{eq:adaptive3}
\end{equation}

Since the neural-network approximation error is bounded, namely,

\begin{equation}
	\|\varepsilon\|
	\le
	\varepsilon_{0},
\end{equation}

the Lyapunov derivative satisfies

\begin{equation}
	\dot{V}
	\le
	-
	\frac{1}{2}
	\lambda_{\min}(Q)
	\|X\|^{2}
	+
	\|PB\|
	\varepsilon_{0}
	\|X\|.
	\label{eq:adaptive4}
\end{equation}

Therefore, whenever

\begin{equation}
	\|X\|
	\ge
	\frac
	{2\|PB\|\varepsilon_{0}}
	{\lambda_{\min}(Q)},
	\label{eq:adaptive5}
\end{equation}

one has $\dot{V}\le0$. Hence, the closed-loop system is uniformly ultimately bounded.

\begin{remark}
	The adaptive law \eqref{eq:adaptive2} is a pure integrator driven by $h(X)X^{T}PB$, which depends only on the state $X$ and contains no term that penalizes the magnitude of $\hat{W}$ itself. Consequently, \eqref{eq:adaptive2} alone cannot guarantee the boundedness of $\hat{W}$: whenever a persistent disturbance or a nonvanishing approximation error $\varepsilon$ is present, the tracking error $X$ no longer converges to the origin but only to a bounded neighborhood as shown in \eqref{eq:adaptive5}, so the driving term $h(X)X^{T}PB$ generally never vanishes. If this residual excitation does not average out to zero, $\hat{W}$ may keep accumulating in a fixed direction indefinitely, leading to unbounded parameter estimates even though the tracking error itself remains bounded. This is the well-known parameter drift phenomenon in adaptive control.
\end{remark}

To suppress parameter drift and endow the adaptive law with an explicit self-decaying mechanism against persistent disturbances and approximation errors, the following $\sigma$-modification is introduced:

\begin{equation}
	\dot{\hat{W}}
	=
	\gamma
	h(X)
	X^{T}
	PB
	-
	\sigma
	\gamma
	\|X\|
	\hat{W},
	\label{eq:sigma1}
\end{equation}

where $\sigma>0$ denotes the leakage coefficient.

Substituting \eqref{eq:sigma1} into \eqref{eq:adaptive1} gives

\begin{equation}
	\dot{V}
	=
	-
	\frac{1}{2}
	X^{T}QX
	+
	X^{T}PB\varepsilon
	+
	\sigma
	\|X\|
	\operatorname{tr}
	\left(
	\tilde{W}^{T}
	\hat{W}
	\right).
	\label{eq:sigma2}
\end{equation}

\subsection{Uniform Ultimate Boundedness Analysis}

To analyze the influence of the $\sigma$-modification term, the following identity is first established:

\begin{equation}
	\operatorname{tr}
	\left(
	\tilde{W}^{T}\hat{W}
	\right)
	=
	\operatorname{tr}
	\left(
	\tilde{W}^{T}W^{*}
	\right)
	-
	\operatorname{tr}
	\left(
	\tilde{W}^{T}\tilde{W}
	\right).
	\label{eq:uub1}
\end{equation}

Applying the Cauchy--Schwarz inequality gives

\begin{equation}
	\operatorname{tr}
	\left(
	\tilde{W}^{T}W^{*}
	\right)
	\le
	\|\tilde{W}\|_{F}
	\|W^{*}\|_{F}
	\le
	\|\tilde{W}\|_{F}
	W_{\max},
	\label{eq:uub2}
\end{equation}

where $W_{\max}$ is a positive constant satisfying

\[
\|W^{*}\|_{F}\le W_{\max}.
\]
Substituting \eqref{eq:uub1} and \eqref{eq:uub2} into \eqref{eq:sigma2} yields
\begin{equation}
	\begin{aligned}
		\dot{V} \le & -\frac{1}{2} \lambda_{\min}(Q) \|X\|^{2} + \|PB\| \varepsilon_{0} \|X\| \\
		& + \sigma \|X\| \|\tilde{W}\|_{F} W_{\max} - \sigma \|X\| \|\tilde{W}\|_{F}^{2}.
	\end{aligned}
	\label{eq:uub3}
\end{equation}

Factoring out $\|X\|$ gives
\begin{equation}
	\begin{aligned}
		\dot{V} \le -\|X\| \Big[ & \frac{1}{2} \lambda_{\min}(Q) \|X\| + \sigma \|\tilde{W}\|_{F}^{2} \\
		& - \sigma W_{\max} \|\tilde{W}\|_{F} - \|PB\| \varepsilon_{0} \Big].
	\end{aligned}
	\label{eq:uub4}
\end{equation}

Completing the square gives
\begin{equation}
	\begin{aligned}
		\sigma \|\tilde{W}\|_{F}^{2} - \sigma W_{\max} \|\tilde{W}\|_{F} = & \, \sigma \left( \|\tilde{W}\|_{F} - \frac{1}{2} W_{\max} \right)^{2} \\
		& - \frac{\sigma}{4} W_{\max}^{2}.
	\end{aligned}
	\label{eq:uub5}
\end{equation}

Substituting \eqref{eq:uub5} into \eqref{eq:uub4} yields
\begin{equation}
	\begin{aligned}
		\dot{V} \le -\|X\| \bigg[ & \frac{1}{2} \lambda_{\min}(Q) \|X\| + \sigma \left( \|\tilde{W}\|_{F} - \frac{1}{2} W_{\max} \right)^{2} \\
		& - \left( \frac{\sigma}{4} W_{\max}^{2} + \|PB\| \varepsilon_{0} \right) \bigg].
	\end{aligned}
	\label{eq:uub6}
\end{equation}

Therefore,

\begin{equation}
	\|X\|
	>
	\frac{2}{\lambda_{\min}(Q)}
	\left(
	\frac{\sigma}{4}
	W_{\max}^{2}
	+
	\|PB\|
	\varepsilon_{0}
	\right),
	\label{eq:uub7}
\end{equation}

and

\begin{equation}
	\|\tilde{W}\|_{F}
	>
	\frac12
	W_{\max}
	+
	\sqrt{
		\frac14
		W_{\max}^{2}
		+
		\frac{\|PB\|
			\varepsilon_{0}}
		{\sigma}
	},
	\label{eq:uub8}
\end{equation}

imply that

\[
\dot V<0.
\]

Hence, both the system state $X$ and the neural-network weight estimation error $\tilde{W}$ are uniformly ultimately bounded.

\begin{remark}
	The conditions in \eqref{eq:uub7} and \eqref{eq:uub8} are alternative conditions. The satisfaction of either condition is sufficient to guarantee that $\dot V<0$, which ensures the uniform ultimate boundedness of the closed-loop system.
\end{remark}

Consequently, all closed-loop signals remain bounded, and the tracking error converges to a compact neighborhood of the origin despite the presence of modeling uncertainties and external disturbances.